\section{Numerical simulation}
\label{sec:simulation}

\subsection{Train--track--bridge interaction model}
\label{sec:ttbi}

Vehicle responses are generated with a two-dimensional, vertically coupled
train--track--bridge interaction (TTBI) model implemented in MATLAB and derived
from the TTB-2D/VEqMon2D framework of Cantero
\cite{cantero_ttb2d, cantero2022veqmon}. The model couples (i) an
Euler--Bernoulli multi-span deck on vertical support springs and abutment
rotational springs; (ii) a discrete layered track containing rail, pads,
sleepers, ballast, and sub-ballast; and (iii) five identical planar vehicles
with car-body and two-bogie vertical/pitch degrees of freedom and four
constrained unsprung wheelset masses. Vehicle properties follow the
implemented intercity-train calibration attributed to O'Brien et
al.~\cite{obrien2018vehicle}. Wheel--rail interaction is linear and bilateral.
The directly coupled equations are integrated by average-acceleration
Newmark-$\beta$ ($\gamma=1/2$, $\beta=1/4$) at 1~ms.

\paragraph{Implementation provenance.}
The generator is a modified derivative of the published TTB-2D v1 code line,
based on archived upstream commit \texttt{28d35528} and redistributed under
GPL-3.0. The damage mechanisms, campaign catalogues, serialization, and
qualification gates described here are repository-local additions and must not
be attributed to the upstream author. Likewise, the track-property file cites
Zhai et al.~\cite{zhai2004track} and the vehicle file cites O'Brien et
al.~\cite{obrien2018vehicle}; those citations establish the source files'
provenance, not independent validation of every transferred value or topology.

Two bridge geometries are used (Figure~\ref{fig:model}). \textbf{F40} has two
20~m spans and one scour target at the internal support. \textbf{L99} has four
24.9~m spans (99.6~m total) and targets at the three internal supports. The
adopted $E$, $I$, mass per unit length, and 3\% bridge-damping target originate
from the Fernandes two-span example; their use for L99 is an idealized
geometry/scale stress, not calibration of a 99.6~m field bridge. Every state
must pass the registered broad frequency gate and healthy states must reproduce
the nominal healthy first frequency (approximately 4.13~Hz for F40 and
2.75~Hz for L99) within the authorized acceptance bands.

\begin{figure}[htbp]
\centering
\framebox[0.95\linewidth]{\rule{0pt}{4.2cm}\parbox{0.9\linewidth}{\centering
\footnotesize Figure placeholder --- F40 and L99 decks on vertical support
springs, simplified layered track, five-vehicle train, and eight
\texttt{physical8\_v1} channels on the leading vehicle. Redraw from the
authorized implementation.}}
\caption{Two-dimensional TTBI model, bridge geometries, and candidate modeled
response channels.}
\label{fig:model}
\end{figure}

\subsubsection{Geometry-specific support-aligned meshes}
\label{sec:mesh}

The production meshes are not universally 0.30~m. F40 uses a 0.20~m deck mesh
(three elements per 0.60~m sleeper bay) and a 0.30~m rail mesh (two per bay), so
the internal support is exactly a deck node at 20.0~m. This is a
support-alignment correction to the Fernandes-derived 0.30~m benchmark density:
0.30~m does not divide 20 or 40~m. L99 uses 0.30~m for both deck and rail,
which places all 24.9~m-spaced supports on nodes. Positive support springs are
required to coincide with nodes to roundoff tolerance; they are never snapped
to a nearby grid location.

Production Rayleigh coefficients are recomputed from the first two elastic
modes of each realized bridge state. Mesh-refinement studies independently
recompute the modal references and coefficients on every grid and save the
resulting $\alpha$, $\beta$, and reference modes. A fixed-coefficient arm using
the healthy reference is a sensitivity, not the production closure. Because
the rail-damping matrix uses the corresponding bridge frequencies as its
references, a structural-stiffness intervention in production changes both the
assembled $K$ and Rayleigh-dependent $C$ matrices; it is not evidence that
physical damage itself changes damping.

\subsubsection{Finite rail-domain selection}
\label{sec:rail_clearance}

The finite rail domain was selected by the source-locked 18-case matrix recorded
in \nolinkurl{docs/evidence/paper1_model_form_freeze_20260809.json}. The
production decision identifier is
\nolinkurl{paper1-rail-domain-clearance-c06-v1}.
For each of F40 and L99, the matrix crossed three operating points
$(70~\mathrm{km/h},3\,^{\circ}\mathrm{C})$,
$(80~\mathrm{km/h},18\,^{\circ}\mathrm{C})$, and
$(90~\mathrm{km/h},33\,^{\circ}\mathrm{C})$ with 6, 15, and 30~m clearances.
All 18 cases completed. Both registered comparison tiers --- C06 versus C30
and C15 versus C30 --- passed, so the rule selected 6~m for production. The
verdict is bound to generator source-root SHA-256
\nolinkurl{aa187204cb3f89e24cb8bc894034044bad38b0358f5e0cd586338f84a8418efb}.
Its authorized claim is strictly finite-rail-domain selection for this
source-locked case matrix; it is not general mesh/time-step convergence,
damage-model validation, or physical/field validation.

\subsubsection{Track-property and topology boundary}
\label{sec:track_boundary}

The implemented nominal track is an inherited hybrid planar convention, not a
blanket sum of all Zhai quantities into two rails. Zhai et
al.~\cite{zhai2004track} report their Table~1 quantities per rail seat at
0.545~m spacing. The implementation doubles rail inertia and rail mass and uses
a full sleeper mass, but retains the tabulated pad, ballast, and sub-ballast
terms at one-seat values. It then places those values on a 0.600~m sleeper
lattice without re-evaluating the spacing-dependent $M_b$, $K_b$, and $K_f$
expressions. The production baseline is therefore neither a demonstrated
one-seat line model nor a consistently doubled two-rail equivalent.

The topology is also simplified relative to Zhai: adjacent-ballast shear
couplings $K_w,C_w$ are omitted, and the complete retained ballast mass is
condensed onto the support-aligned deck degree of freedom at every on-bridge
sleeper. This is an explicit inherited model-form choice. The topology arm was
dropped because it would require new solver structure and is not examined in
this paper.

Three parameter-level sensitivity families are prospectively fixed against the
hybrid baseline: (i) consistent one-seat and consistent two-rail scaling;
(ii) recomputation of $M_b$, $K_b$, and $K_f$ at 0.600~m from Zhai's equations;
and (iii) 0.050\% and 0.200\% rail-damping targets around the inherited 0.100\%
target. The 0.100\% target is author-chosen and is not reported by Zhai.
Sensitivity outcomes qualify interpretation but do not silently redefine the
production model after its source freeze.

\begin{table}[htbp]
\centering
\caption{Nominal physical and numerical parameters read from the implemented
baseline. Track entries preserve the hybrid scaling and topology described in
Section~\ref{sec:track_boundary}; they must not be labeled uniformly
``two-rail Zhai values''.}
\label{tab:parameters}
\footnotesize
\setlength{\tabcolsep}{4pt}
\resizebox{\textwidth}{!}{%
\begin{tabular}{llll}
\toprule
Subsystem & Parameter & Value & Units \\
\midrule
\multirow{8}{*}{Bridge}
 & spans (F40 / L99) & $2\times20$ / $4\times24.9$ & m \\
 & Young's modulus $E$ at 15\,$^{\circ}$C & $35\times10^{9}$ & N/m$^2$ \\
 & second moment of area $I$ & 0.33 & m$^4$ \\
 & mass per unit length & 9600 & kg/m \\
 & Rayleigh target, first two elastic modes & 3 & \% \\
 & healthy vertical support stiffness $k_{v,0}$ & $3.44\times10^{8}$ & N/m \\
 & free-rotation abutment baseline & 0 & N$\cdot$m/rad \\
 & healthy first frequency (F40 / L99) & $\approx4.13$ / $\approx2.75$ & Hz \\
\midrule
\multirow{8}{*}{Track}
 & rail bending stiffness $EI$ & $1.325\times10^{7}$ & N$\cdot$m$^2$ \\
 & rail mass (doubled-rail entry) & 121.28 & kg/m \\
 & rail Rayleigh target (inherited) & 0.100 & \% \\
 & pad stiffness / damping (one-seat entry) & $6.5\times10^{7}$ / $7.5\times10^{4}$ & N/m, N$\cdot$s/m \\
 & sleeper spacing / full mass & 0.600 / 251 & m, kg \\
 & ballast mass / stiffness / damping (one-seat entries) & 531.4 / $1.3775\times10^{8}$ / $5.88\times10^{4}$ & kg, N/m, N$\cdot$s/m \\
 & sub-ballast stiffness / damping (one-seat entries) & $7.75\times10^{7}$ / $3.115\times10^{4}$ & N/m, N$\cdot$s/m \\
 & finite rail clearance beyond deck & 6 & m \\
\midrule
\multirow{7}{*}{Vehicle ($\times5$)}
 & car-body mass / pitch inertia & 36852 / 560342 & kg, kg$\cdot$m$^2$ \\
 & bogie mass / pitch inertia & 3910 / 10024 & kg, kg$\cdot$m$^2$ \\
 & wheelset mass & 1407 & kg \\
 & primary suspension stiffness / damping per axle & $5.558\times10^{6}$ / $5.88\times10^{4}$ & N/m, N$\cdot$s/m \\
 & secondary suspension stiffness / damping per bogie & $2.0\times10^{6}$ / $1.2\times10^{5}$ & N/m, N$\cdot$s/m \\
 & bogie-centre / axle spacing & 16 / 2.3 & m \\
 & vehicle length / wheel radius & 22 / 0.46 & m \\
\midrule
\multirow{4}{*}{Solver}
 & integration & Newmark-$\beta$ ($\gamma=1/2$, $\beta=1/4$) & --- \\
 & time step / stored rate & $10^{-3}$ / 1000 & s, Hz \\
 & deck element (F40 / L99) & 0.20 / 0.30 & m \\
 & rail element (F40 / L99) & 0.30 / 0.30 & m \\
\bottomrule
\end{tabular}%
}
\end{table}

\subsection{Response-channel schema}
\label{sec:response_channels}

The leading vehicle supplies the eight candidates in
Table~\ref{tab:channels}, identified collectively as
\texttt{physical8\_v1}. Six are direct responses of modeled car-body/bogie
degrees of freedom. The other two are the total vertical accelerations implied
by the constrained motion of the first two wheelsets,
\begin{equation}
  \ddot z_w = u_{,tt} + 2v\,u_{,xt} + v^2 u_{,xx} + \ddot h_w,
  \label{eq:wheelset_proxy}
\end{equation}
where $u$ is the rail displacement field evaluated along the moving contact
coordinate, $v$ is passage speed, and $h_w$ is the prescribed profile/wheel
path. These two rows are idealized, model-predicted
\emph{constrained-wheelset acceleration proxies}. They are not independent
wheelset degrees of freedom, not the Eulerian rail-field acceleration
$u_{,tt}$ alone, and not validated measurements from axle-box sensors.
The separate stored \texttt{acc\_under} rows remain virtual rail-field
diagnostics and never enter the eight-channel screen.

\begin{table}[htbp]
\centering
\caption{Candidate \texttt{physical8\_v1} responses of the leading vehicle.}
\label{tab:channels}
\small
\begin{tabular}{cll}
\toprule
Index & Response channel & Observation status \\
\midrule
0 & car-body vertical acceleration & modeled secondary-suspended DOF \\
1 & front-bogie vertical acceleration & modeled primary-suspended DOF \\
2 & rear-bogie vertical acceleration & modeled primary-suspended DOF \\
3 & wheelset-1 constrained vertical acceleration & idealized total-motion proxy \\
4 & wheelset-2 constrained vertical acceleration & idealized total-motion proxy \\
5 & car-body pitch angular velocity & modeled secondary-suspended DOF \\
6 & front-bogie pitch angular velocity & modeled primary-suspended DOF \\
7 & rear-bogie pitch angular velocity & modeled primary-suspended DOF \\
\bottomrule
\end{tabular}
\end{table}

The proxy label is deliberate. Equation~\eqref{eq:wheelset_proxy} omits a
hardware observation chain, local axle-box mounting dynamics, sensor bandwidth
and filtering, and contact nonlinearity. Accordingly, the channel screen ranks
modeled responses and does not determine a transducer count or procurement
choice.

\subsubsection{Contact-model scope and qualification}
\label{sec:contact}

The bilateral linear contact model does not simulate separation and re-contact.
Wheel flats and polygonization are disabled in all four Paper-1 blocks. Every
retained passage nevertheless carries finite/contact diagnostics and fails
generation if it violates the registered admissibility contract. Before
dispatch, the contact-closure study evaluates all 420 passages of the four
marked qualification datasets at 1.0, 0.5, and 0.25~ms, applies the registered
0/12/24~kN tension gates and 0.002 path-fraction gate, and checks waveform and
quantity-of-interest convergence. This is an engineering consistency gate for
the finite design, not validation of bilateral contact physics.

\subsection{Active bridge mechanisms}
\label{sec:mechanisms}

\subsubsection{Scour as support-stiffness loss}
\label{sec:scour_mech}

At each target support, scour is represented as
\begin{equation}
  k_v(d)=(1-d)k_{v,0}, \qquad 0\le d\le0.60,
  \label{eq:scour}
\end{equation}
with $k_{v,0}=3.44\times10^8$~N/m. The regression label is $100d$,
modeled support-stiffness loss in percentage points. This follows the
drive-by-scour convention~\cite{fernandes2024driveby,
fernandes2025early, prendergast2013investigation}, but is not scour depth,
embedment loss, or eroded volume.

\subsubsection{Nominal bearing fixity}
\label{sec:bearing_mech}

F40-M and L99-M activate a rotational spring at each abutment. The sampled
geometry-normalized coordinate is
\begin{equation}
 \varphi=\frac{k_r}{k_r+4E_{15}I/L_{\mathrm{end}}}
 \quad\Longleftrightarrow\quad
 k_r=\frac{\varphi}{1-\varphi}\frac{4E_{15}I}{L_{\mathrm{end}}},
 \qquad \varphi\in[0,0.95].
 \label{eq:fixity}
\end{equation}
The stiffness is computed once using the fixed 15\,$^{\circ}$C modulus and
remains constant while the deck modulus varies between passages. The two
fixity coordinates are additional M-block regression targets. They are nominal
free-to-near-fixed design variables, not calibrated bearing-damage percentages
or physical seizure thresholds.

\subsubsection{Local flexural-rigidity-loss surrogate}
\label{sec:crack_mech}

F40-M and L99-M reduce $EI$ uniformly over the single production finite element
containing the sampled damage location, with author-chosen severity
$\mathcal{U}(0.05,0.30)$. Locations use an author-chosen 4:1 design weighting
toward hogging regions inside the registered support-zone and bridge-length
bounds. No fitting dataset supports those occurrence or severity laws. The
damaged length is consequently one mesh-locked element --- 0.20~m on F40 and
0.30~m on L99 --- rather than a common physical crack length. This is a local
damaged-element benchmark, not the tapered crack-depth model of Sinha et
al.~\cite{sinha2002simplified}; its severity is never interpreted as crack
depth. Crack status is logged but not predicted.

\subsection{Fixed profile, operational variability, and deferred mechanisms}
\label{sec:opvar}

Every state and passage in all four blocks uses one generated realization of
the implemented FRA-class-4 vertical-profile spectrum (registered phase seed
20260728). No measured profile, per-state phase variation, or per-passage
profile jitter is used. The fixed choice enables exact paired comparisons but
removes profile variability from the production EOV distribution.

Operational conditions do vary. Each state's 50 passages use a Latin
hypercube over speed 70--90~km/h and temperature 3--33~$^{\circ}$C, rounded to
integer values. Temperature changes the deck modulus through the registered
author-chosen law
\[
 E(T)=E_{15}\,[1-0.003(T-15)].
\]
For each passage and vehicle, car-body mass has a 10\% Gaussian coefficient of
variation and primary/secondary suspension stiffnesses each have 5\%. These
envelopes and laws were not fitted to a route, climate, or fleet population.

Track-layer damage (ballast patches, hanging sleepers, and pad failure), wheel
polygonization, suspension damage, and alternative rail-profile phases are
disabled by all four stage contracts. The corresponding code remains in the
single source tree for future work; it contributes no expressed physics or
random variability to Paper~1.

\subsection{State catalogues and common random numbers}
\label{sec:states}

Table~\ref{tab:families} gives the exact state inventories. F40-S is a dense
61-severity grid (0--60\% in one-percentage-point increments) with five
semantic replicas per severity. F40-M uses five controlled families. Its five
healthy replicas and five replicas at each of 12, 24, 36, 48, and 60\% scour
form the exact 30-state subset shared with F40-S. L99-S and L99-M share their
complete five-family, 475-state latent design; bearing/crack coordinates are
dormant in L99-S rather than being removed or redrawn.

\begin{table}[htbp]
\centering
\caption{Registered state families. Every state has 50 operational passages.}
\label{tab:families}
\small
\begin{tabular}{lrrrp{5.0cm}}
\toprule
Family & F40-S & F40-M & L99-S/M & Role \\
\midrule
\texttt{target\_healthy} & 5 & 50 & 50 & zero-scour diagnostic \\
\texttt{scour\_only} & 300 & 25 & 75 & controlled support/severity anchors \\
\texttt{bearing\_only} & 0 & 50 & 50 & bearing/leakage anchors; dormant in L99-S \\
\texttt{nuisance\_only} & 0 & 50 & 50 & local-$EI$ anchors; dormant in L99-S \\
\texttt{joint} & 0 & 250 & 250 & multivariate latent design \\
\midrule
Total & 305 & 425 & 475 & \\
\bottomrule
\end{tabular}
\end{table}

Each state has a collision-checked semantic \texttt{StateUID}. Named substreams
derive operational, crack, profile, track, and passage identities from that UID
rather than a mutable row number or parallel-execution order. Only streams for
expressed Paper-1 mechanisms influence the solve; disabled track/wheel
namespaces remain inert. States are the split and resampling clusters, and
their 50 passages are correlated repeated observations, not independent field
samples.

The generator stores raw noise-free histories, semantic identities, latent
design coordinates, operational draws, profile and model-form identifiers,
modal/Rayleigh records, contact diagnostics, crop metadata, and source-bound
fingerprints. Observation noise is added only by the authenticated learning
pipeline (Section~\ref{sec:preprocessing}).
